\documentclass[11pt,a4paper]{article}
\usepackage[margin=1in]{geometry}
\usepackage{hyperref}
\usepackage{enumitem}
\usepackage[utf8]{inputenc}
\usepackage[T1]{fontenc}
\usepackage[portuguese]{babel}

\begin{document}

\begin{center}
    {\Huge \textbf{Prabhavit Gupta}} \\
    \vspace{2mm}
    Via Zanetti 10, 30175, Veneza, Itália $\vert$ +39 351 610 7807 $\vert$ \href{mailto:prabhavitg@gmail.com}{prabhavitg@gmail.com} \\
    \href{https://www.linkedin.com/in/prabhavit/}{linkedin.com/in/prabhavit/} $\vert$ \href{https://github.com/prab-gupta}{github.com/prab-gupta}
\end{center}

\vspace{5mm}
\section*{Resumo Profissional}
Desenvolvedor Full Stack Sênior com ampla experiência na criação de soluções de software de ponta a ponta e escaláveis. Especialista em ecossistemas modernos utilizando React, Angular, Node.js e NestJS. Forte capacidade em desenhar APIs RESTful robustas, gerenciar infraestrutura em nuvem na AWS e manter pipelines de CI/CD eficientes com o Azure DevOps. Experiente com bancos de dados relacionais e não-relacionais, sempre focado na entrega de soluções de alta qualidade em ambientes colaborativos e dinâmicos.

\section*{Habilidades Técnicas}
\begin{itemize}[leftmargin=*]
    \item \textbf{Frontend:} React, Angular, TypeScript, JavaScript, HTML5, CSS3, RxJS
    \item \textbf{Backend:} Node.js, NestJS, C\#, .NET Core, Python
    \item \textbf{Bancos de Dados:} SQL Server, PostgreSQL, MongoDB, Redis
    \item \textbf{Nuvem \& DevOps:} AWS, Docker, Azure DevOps, CI/CD, Git
    \item \textbf{Ferramentas \& Metodologias:} Postman, APIs RESTful, TDD, Agile
\end{itemize}

\section*{Experiência Profissional}
\noindent
\textbf{Desenvolvedor Full Stack Sênior} \hfill 2021 -- Presente \\
\textit{Autônomo / Contrato} \hfill Remoto
\begin{itemize}[leftmargin=*]
    \item Desenhou e implementou aplicações web robustas de ponta a ponta, utilizando Node.js e NestJS no backend, integrados a interfaces dinâmicas em React e Angular.
    \item Construiu, testou e otimizou APIs RESTful seguras e escaláveis, resultando na redução do tempo de resposta e melhoria do desempenho do sistema.
    \item Gerenciou a infraestrutura em nuvem utilizando serviços AWS e aplicações conteinerizadas com Docker, assegurando estabilidade e alta disponibilidade.
    \item Utilizou Azure DevOps para o gerenciamento ágil de projetos, automação de processos e manutenção de pipelines de CI/CD eficientes.
    \item Modelou esquemas de banco de dados e otimizou consultas complexas tanto em sistemas relacionais (PostgreSQL) quanto não-relacionais (MongoDB, Redis).
\end{itemize}

\vspace{2mm}
\noindent
\textbf{Engenheiro de Software} \hfill 2018 -- 2021 \\
\textit{Tech Solutions Inc.} \hfill Milão, Itália (Híbrido)
\begin{itemize}[leftmargin=*]
    \item Contribuiu para o desenvolvimento de sistemas de software corporativos utilizando frameworks Full Stack modernos.
    \item Manteve e aprimorou pipelines de integração contínua (CI/CD) com Docker, diminuindo os tempos de implantação e entrega de software.
    \item Aplicou práticas de código limpo e desenvolvimento orientado a testes (TDD), garantindo a cobertura abrangente de testes unitários.
\end{itemize}

\section*{Educação}
\noindent
\textbf{Mestrado em Ciência da Computação} \hfill 2018 \\
\textit{Universidade de Veneza} \hfill Veneza, Itália

\end{document}
